% Generated draft: requires mathematical and source review.
\documentclass[11pt]{article}
\usepackage[margin=25mm]{geometry}
\usepackage{fontspec}
\setmainfont{Latin Modern Roman}
\usepackage{amsmath,amssymb,mathtools,unicode-math}
\setmathfont{Latin Modern Math}
\usepackage{enumitem,needspace,xurl,hyperref}
\setlength{\parindent}{0pt}
\setlength{\parskip}{5pt}
\setlength{\emergencystretch}{3em}
\newcommand{\N}{\mathbb{N}}
\newcommand{\Z}{\mathbb{Z}}
\newcommand{\Q}{\mathbb{Q}}
\newcommand{\R}{\mathbb{R}}
\newcommand{\C}{\mathbb{C}}
\newcommand{\F}{\mathbb{F}}
\newcommand{\A}{\mathbb{A}}
\newcommand{\PP}{\mathbb{P}}
\newcommand{\E}{\mathbb{E}}
\newcommand{\eps}{\varepsilon}
\newcommand{\dd}{\,\mathrm{d}}
\begin{document}
\section*{UnsolvedMath ID: 11000099}
\textbf{Source problem code:} catalogue entry
\subsection{Definitions and mathematical statement}
Let $S$ be the surface in the source and $\operatorname{Mod}_S$ its mapping class group. In the usual multitwist convention, choose finitely many pairwise disjoint, nonisotopic essential curves $\alpha_1,\ldots,\alpha_n$ and signs $\epsilon_i\in\{1,-1\}$, and put $t=t_{\alpha_1}^{\epsilon_1}\cdots t_{\alpha_n}^{\epsilon_n}$. For $f\in\operatorname{Mod}_S$, assume the family
\[
 A=\{f^m(\alpha_i):1\leq i\leq n,\ m\in\Z\}
\]
fills $S$: every essential nonperipheral simple closed curve $\gamma$ has positive geometric intersection number with some member of $A$. The source says ``nontrivial circle'' and does not explicitly state these standard curve conventions.

Question I asks whether there is a positive integer $N_S$, depending only on $S$, such that for every such $f,t$ the exceptional set
\[
 E(f,t)=\{j\in\Z:t^j\circ f\text{ is not pseudo-Anosov}\}
\]
is contained in a block of $N_S$ consecutive integers. Precisely, there should exist $a=a(f,t)\in\Z$ with
\[
 E(f,t)\subset\{a,a+1,\ldots,a+N_S-1\}.
\]
This includes finiteness and a bound on the span of exceptions. The source separately asks the weaker question whether $\#E(f,t)\leq N_S$, without requiring the exceptional indices to be consecutive or close together.

The source attributes the single-Dehn-twist case with $N_S=7$ to A. Fathi [Fath]. It says the original motivation, proving omitted Conjecture 4 in certain cases, is now obsolete (with a qualification about complete obsolescence). It also cites an application [I2] giving mapping-class-group rank one in an unspecified sense and [I7], Section 9.4. These contextual statements do not supply a proof for multitwists.
\subsection{Short English statement}
For a multitwist whose curves have filling orbit under $f$, are all non-pseudo-Anosov members of $t^j f$ confined to a uniformly bounded block of exponents? A weaker question bounds only the number of exceptions.
\subsection{Sources}
\begin{itemize}
\item Local registry: \url{data/unsolvedmath/problems.json}, permanent UnsolvedMath ID $11000099$, exact source code \texttt{catalogue entry}.
\item Farb, \emph{Problems on Mapping Class Groups and Related Topics}, Question I, PDF page 83: \url{https://www.math.uchicago.edu/~farb/papers/mcgbook.pdf}. Supplied access date: 2026-07-29; extraction: pdf-font-structure; repeated local numbering disambiguated by page among 329 headings.
\item Supplied citations: A. Fathi [Fath] for the seven-consecutive-exponent result, [I2], [I7] Section 9.4, and Conjecture 4. Full bibliography entries and the conjecture are absent. The registry literature summary is explicitly truncated.
\end{itemize}
\end{document}
